% Part: lambda-calculus
% Chapter: introduction
% Section: primitive-recursion

\documentclass[../../../include/open-logic-section]{subfiles}

\begin{document}

\olfileid[pt]{lam}{int}{pr}
\olsection{As funções \usetoken{S}{lambda definable} são fechadas sob recursão primitiva}

Quando chegamos à recursão primitiva, finalmente precisamos trabalhar um
pouco. Teremos de avançar por etapas. Como antes, supondo que já tenhamos
termos $G$ e $H$ que !!{lambda define}m as funções $g$ e~$h$,
respectivamente, queremos um termo $H$ que !!{lambda define}s a função~$f$
definida por
\begin{align*}
f(0, \vec z) & = g(\vec z) \\
f(x+1, \vec z) & = h(z, f(x,\vec z), \vec z).
\end{align*}
Assim, em geral, dados os termos lambda $G'$ e~$H'$, basta encontrar um
termo~$F$ tal que
\begin{align*}
F(\num{0}, \vec z) & \equiv G(\vec z) \\
F(\overline{n+1}, \vec z) & \equiv H(\num{n}, F(\num{n}, \vec z), \vec z)
\end{align*}
para todo número natural~$n$; o fato de $G'$ e~$H'$ !!{lambda define}rem
$g$ e~$h$ significa que, sempre que inserirmos numerais $\num{\vec m}$ no
lugar de $\vec z$, $F(\num{n+1}, \num{\vec m})$ se normalizará para a
resposta correta.

Para isso, porém, basta encontrar um termo~$F$ que satisfaça
\begin{align*}
F(\num 0) & \equiv G \\
F(\overline {n+1}) & \equiv H(\num{n},F(\num{n}))
\intertext{para todo número natural $n$, onde}
G & = \lambd[\vec z][G'(\vec z)] \text{ e}\\
H(u,v) & = \lambd[\vec z][H'(u,v(u,\vec z),\vec z)].
\end{align*}
Em outras palavras, com artifícios lambda podemos evitar a preocupação com
os parâmetros adicionais $\vec z$---eles são simplesmente absorvidos pela
notação lambda.

Antes de definir o termo $F$, precisamos de um mecanismo para lidar com
pares ordenados. Ele é fornecido pelo próximo lema.

\begin{lem}
Existe um termo lambda $D$ tal que, para cada par de termos lambda $M$ e~$N$,
$D(M,N)(\num{0}) \red M$ e $D(M,N)(\num{1}) \red N$.
\end{lem}

\begin{proof}
Primeiro, defina o termo lambda $K$ por
\[
K(y) = \lambd[x][y].
\]
Em outras palavras, $K$ é o termo $\lambd[y][\lambd[x][y]]$. Vendo de outra
maneira, para todo $M$, $K(M)$ é uma função constante que devolve~$M$ para
qualquer entrada.

Agora defina $D(x,y,z)$ por $D(x,y,z) = z (K(y))x$. Então temos
\begin{align*}
D(M,N,\num 0) & \red \num 0 (K(N)) M \red M \text{ e}\\
D(M,N,\num 1) & \red \num 1 (K(N)) M \red K(N) M \red N,
\end{align*}
como requerido.
\end{proof}

A ideia é que $D(M,N)$ representa o par $\tuple{M, N}$ e, se supomos que
$P$ representa tal par, $P(\num 0)$ e $P(\num 1)$ representam as projeções
esquerda e direita, $(P)_0$ e $(P)_1$. Usaremos estas últimas notações.

\begin{lem}
As funções !!{lambda definable} são fechadas sob recursão primitiva.
\end{lem}

\begin{proof}
Precisamos mostrar que, dados termos quaisquer $G$ e $H$, podemos encontrar
um termo $F$ tal que
\begin{align*}
F(\num{0}) & \equiv G \\
F(\num{n+1}) & \equiv H(\num{n}, F(\num{n}))
\end{align*}
para todo número natural~$n$. A ideia, em linhas gerais, é calcular
sequências de \emph{pares}
\[
\tuple{\num 0, F(\num 0)}, \tuple{\num 1, F(\num 1)}, \dots,
\]
usando numerais como iteradores. Observe que o primeiro par é simplesmente
$\tuple{\num 0, G}$. Dado um par $\tuple{\num{n}, F(\num{n})}$, o par
seguinte, $\tuple{\num{n+1}, F(\num{n+1})}$, deve ser equivalente a
$\tuple{\num{n+1}, H(\num{n}, F(\num{n}))}$. Projetaremos um termo
lambda~$T$ que realiza essa transição em um passo.

Os detalhes são os seguintes. Defina $T(u)$ por
\[
T(u) = \tuple{S((u)_0), H((u)_0,(u)_1)}.
\]
Agora é fácil verificar que, para qualquer número~$n$,
\[
T(\tuple{\num{n}, M}) \red \tuple{\num{n+1}, H(\num{n}, M)}.
\]
Como sugerido acima, dados $G$ e $H$, defina~$F(u)$ por
\[
F(u) = (u(T,\tuple{\num 0, G}))_1.
\]
Em outras palavras, com entrada~$\num{n}$, $F$ itera $T$ $n$ vezes sobre
$\tuple{\num 0, G}$ e depois devolve a segunda componente. Para começar,
temos
\begin{enumerate}
\item $\num{0} (T, \tuple{\num 0, G}) \equiv \tuple{\num{0}, G}$
\item $F(\num{0}) \equiv G$
\end{enumerate}
Por indução em~$n$, podemos mostrar que, para cada número natural, vale o
seguinte:
\begin{enumerate}
\item $\num{n+1}(T, \tuple{\num{0}, G}) \equiv \tuple{\num{n+1}, F(\num{n+1})}$
\item $F(\num{n+1}) \equiv H(\num{n}, F(\num{n}))$
\end{enumerate}
Para a segunda cláusula, temos
\begin{align*}
F(\num{n+1}) & \red (\num{n+1}(T, \tuple{\num 0, G}))_1 \\
& \equiv (T(\num{n} (T, \tuple{\num 0, G})))_1 \\
& \equiv (T(\tuple{\num{n}, F(\num{n})}))_1 \\
& \equiv (\tuple{\num{n+1}, H(\num{n}, F(\num{n}))})_1 \\
& \equiv H(\num{n}, F(\num{n})).
\end{align*}
Aqui usamos a hipótese de indução na penúltima linha. Para a primeira
cláusula, temos
\begin{align*}
\num{n+1} (T, \tuple{\num 0, G}) &
\equiv T(\num{n} (T, \tuple{\num 0, G})) \\
& \equiv T( \tuple{\num{n}, F(\num{n})}) \\
& \equiv \tuple{\num{n+1}, H(\num{n}, F(\num{n}))} \\
& \equiv \tuple{\num{n+1}, F(\num{n+1})}.
\end{align*}
Aqui usamos a segunda cláusula na última linha. Portanto, mostramos
$F(\num 0) \equiv G$ e, para todo $n$, $F(\num {n+1}) \equiv H(\num
n, F(\num{n}))$, exatamente o que precisávamos.
\end{proof}

\end{document}


